\section{Introduction}

Immune tissues are composed of multiple cell populations, each with a distinct transcriptional program that can independently modulate disease risk, progression and response to therapy. Cell-type-specific gene expression has been shown to discriminate disease states more effectively than mixed-cell measurements in a range of immune-mediated diseases, and signatures such as CD8 T cell exhaustion predict clinical course across autoimmunity and infection \cite{mckinney2015tcell}. Transcriptomic studies of juvenile idiopathic arthritis in particular have linked the activation state of specific immune subsets to treatment response and disease severity \cite{throm2018juvenile}. Accessing cell-type resolution, however, typically requires flow sorting followed by RNA extraction, or single-cell RNA sequencing (scRNA-seq), both of which are labour- and cost-intensive at the sample sizes needed for well-powered cohort studies \cite{oelen2022singlecell,monaco2019rnaseq}. As a consequence, most large immune disease cohorts are profiled as bulk mixtures, most often peripheral blood mononuclear cells (PBMC), in which cell-type-specific signals are diluted.

Computational deconvolution has emerged as a pragmatic alternative \cite{avilacobos2020benchmarking}. A standard formulation models an observed bulk expression vector $m$ for $n$ genes as $m = H f$, where $H$ is an $n \times c$ matrix of cell-type-specific expression profiles and $f$ is a vector of $c$ cell-type fractions \cite{shenorr2010cssam}. The majority of methods developed to date focus on estimating $f$ given $H$: EPIC \cite{racle2017epic} and quanTIseq \cite{finotello2019quantiseq} use constrained regression, FARDEEP \cite{hao2019fardeep} applies adaptive least-trimmed squares, and CIBERSORT \cite{newman2015robust} applies linear $\nu$-support vector regression. These methods differ primarily in how they regularise the regression against a reference signature matrix, and they have been extensively benchmarked in synthetic mixtures with known ground truth \cite{avilacobos2020benchmarking}.

Imputing sample-level cell-type expression, the three-way tensor $G$ of size $n \times c \times k$ for $k$ samples, is a substantially harder problem and is required for analyses that link expression to quantitative traits or that stratify patients. Only a handful of methods address it directly. CIBERSORTx \cite{newman2019cibersortx} extends CIBERSORT with a high-resolution module that iteratively refines per-gene estimates under a non-negative least squares constraint. MIND \cite{wang2020mind} and its Bayesian extension bMIND \cite{wang2024bmind} use empirical Bayes and Markov chain Monte Carlo sampling respectively to borrow information across multiple tissue measurements per subject. swCAM \cite{chen2022swcam} is built on the debCAM convex-analysis-of-mixtures framework \cite{chen2020debcam,wang2016undo} and uses low-rank matrix factorisation to recover sample-specific variation. Each of these methods carries restrictive assumptions: MIND requires multiple bulk measurements per subject, bMIND and swCAM expect either flow-measured fractions or a high-quality signature matrix, and all of them rely in some way on the accuracy of a preliminary fraction estimate.

Penalised multi-response regression, by contrast, imposes no such structural assumption. Given paired PBMC and sorted cell RNA-seq data from training subjects, one can directly learn a map from PBMC expression to sorted-cell expression for every gene jointly. Multivariate LASSO and ridge regression, implemented efficiently in glmnet \cite{friedman2010glmnet,tibshirani1996lasso}, exploit shared predictor-response structure so that regression coefficients can be shrunk to zero across groups of correlated target genes. This cross-domain machine-learning approach has been widely used in genomics for phenotype prediction and in eQTL mapping, but to our knowledge has never been systematically compared to domain-specific deconvolution methods for the sample-level cell-type expression imputation problem.

In this work we benchmark three domain-specific methods (CIBERSORTx with inbuilt and custom signatures, bMIND with flow fractions, and swCAM with flow fractions) against multi-response LASSO and ridge regression. We use two complementary datasets. The first is a real-world paired cohort of 158 subjects from the CLUSTER Consortium with matched PBMC and sorted CD4, CD8, CD14 and CD19 RNA-seq, processed through a harmonised STAR \cite{dobin2013star} and featureCounts \cite{liao2014featurecounts} pipeline with Combat-seq batch correction \cite{zhang2020combatseq} and edgeR low-count filtering \cite{robinson2010edger}. The second is a family of pseudobulk datasets synthesised from the eQTLgen 1M-scBloodNL scRNA-seq resource \cite{oelen2022singlecell} using the Jaakkola and Elo construction \cite{jaakkola2017comparison}, which let us vary training size and batch-correction strategy systematically.

Beyond the standard correlation and RMSE metrics, we introduce a differential-gene-expression (DGE) recovery metric that directly measures the usefulness of imputed expression for downstream hypothesis testing. For each simulated phenotype we run limma \cite{ritchie2015limma} in parallel on observed and imputed data, threshold the observed analysis at FDR 0.05, vary the threshold in the imputed analysis from 0 to 1 in steps of 0.05, and compute a receiver-operating-characteristic AUC using pROC \cite{robin2011proc}. Our contributions are: (i) a multi-response LASSO/ridge imputation framework that clusters correlated target genes into chunks (using dynamic hierarchical clustering \cite{langfelder2008dynamictreecut}) so that regression coefficients are shared across correlated outputs; (ii) a task-aligned DGE-recovery metric that exposes performance differences correlation and RMSE cannot see; and (iii) a head-to-head evaluation across 24 pseudobulk datasets and four CLUSTER phenotype scenarios showing that penalised regression dominates domain-specific methods on DGE recovery, at the cost of substantially higher compute.

\section{Related work}

\paragraph{Fraction deconvolution.} The first generation of deconvolution methods focused on estimating cell-type fractions $f$ from bulk mixtures given a reference signature matrix. CIBERSORT \cite{newman2015robust} used $\nu$-support vector regression to gain robustness against collinear signatures and unknown mixture components. EPIC \cite{racle2017epic} and quanTIseq \cite{finotello2019quantiseq} used constrained least-squares regression, and FARDEEP \cite{hao2019fardeep} added adaptive least-trimmed-squares outlier rejection. Fully unsupervised methods such as debCAM \cite{chen2020debcam} and the related UNDO framework \cite{wang2016undo} estimate both signature and fractions jointly by convex analysis of mixtures. Benchmarks across synthetic scRNA-seq-derived mixtures consistently show that fraction estimation is accurate for cell types with distinct signatures but remains challenging for closely related subsets such as CD4 and CD8 T cells \cite{avilacobos2020benchmarking}.

\paragraph{Sample-level cell-type expression.} CIBERSORTx \cite{newman2019cibersortx} extends CIBERSORT with high-resolution imputation by iteratively solving bootstrapped NNLS for each gene, conditional on an earlier fraction estimate. CellR adopts a negative-binomial model and simulated annealing but requires pre-defined scRNA-seq clusters. MIND \cite{wang2020mind} uses an expectation-maximisation algorithm that demands multiple bulk measurements per subject. bMIND \cite{wang2024bmind} replaces this requirement with a Bayesian mixed-effects model that estimates per-subject expression through MCMC. swCAM \cite{chen2022swcam} augments debCAM with a low-rank matrix-factorisation term under nuclear-norm and $\ell_{2,1}$ penalties, recovering per-sample variation from unsupervised mixtures. All of these methods depend either on external fractions or on a signature matrix, and they operate gene-by-gene on the response side, losing the correlation structure between genes that penalised multi-response regression can exploit.

\paragraph{Benchmarks and downstream metrics.} Most existing deconvolution benchmarks use correlation and RMSE per subject as their primary metrics \cite{avilacobos2020benchmarking,monaco2019rnaseq}. As we show below, these metrics are dominated by between-cell-type variance and are only weakly discriminative for expression imputation. Task-specific metrics that quantify recovery of differential expression, as originally proposed in csSAM-style analyses \cite{shenorr2010cssam} and adopted by the bMIND and swCAM evaluations, track downstream utility more faithfully. Our DGE-recovery AUC formalises this intuition into a single scalar metric suitable for side-by-side method comparison.

\paragraph{This work in context.} By treating imputation as a regression problem with paired training data rather than as a constrained inverse problem, we sidestep the fraction-estimation step entirely and inherit all of the machinery of penalised regression: cross-validated regularisation paths, multi-response grouping, and efficient fitting in glmnet \cite{friedman2010glmnet}. The trade-off is explicit: when paired PBMC and sorted training data exist, LASSO and ridge dominate; when they do not, the deconvolution-based methods remain necessary. The CLUSTER cohort studied here, with its four sorted cell types measured alongside PBMC in 80 training subjects, is precisely the regime in which cross-domain machine learning can be expected to win.
